\documentclass[10pt,a4paper]{article}

\usepackage{epsfig}
\usepackage{amsmath}
\usepackage{graphicx}
\usepackage{float}
\usepackage{subfig}
\usepackage{vmargin}
\usepackage{mathrsfs}
\usepackage{mathbbol}
\usepackage[round]{natbib}
\usepackage[parfill]{parskip}

\usepackage{fancyhdr}
\pagestyle{fancy}
\rhead{\includegraphics[width=3.5cm]{../juniper/juniper-light-bg.png}}
\lhead{\includegraphics[width=2.5cm]{../juniper/Manchester.pdf}}

% Here enter any other preamble

\newcommand{\be}{\begin{equation}}
\newcommand{\ee}{\end{equation}}
\newcommand{\ba}{\begin{equation} \begin{aligned}}
\newcommand{\ea}{\end{aligned} \end{equation}}

%\newtheorem{mydef}[equation]{Definition}
\newtheorem{mydef}{Definition}

\title{\sc Importation of Variants and Border Control}
\author{Thomas House \and Jacob Curran Sebastian \and University of Manchester
COVID-19 Modelling Group \and JUNIPER Consortium}
\date{Report for SPI-M 24 March 2021}

\begin{document}

\maketitle
\thispagestyle{fancy}
\section{Summary of key results}

\begin{itemize}
\item We are currently seeing just under a couple of vaccine escape VoC infections
detected per day.
\item Even if this rate continues, it would require stopping over 99\% of new
introductions to have a reasonable chance of stopping introduction before
the 100 or so days required to produce updated vaccine.
\item In the (more realistic) scenario of exponentially growing importations,
success rates need to be over 99.9\% successful to achieve similar results.
\item Even for successfully controlled introductions, the time to extinction
may mean that the UK would almost never be in a position of `zero variants'.
\end{itemize}

\section{Analysis of VoC data}

Data were taken from the \texttt{VOC202012\_02\_linelist\_20210319} file on the
DSTL server. This was separated into exposure types, and the `Traveller' stream
was analysed using the methods from \citet{House:2018}. These indicate
introduction of new infections at a rate of approximately 1.7 per day with no
evidence of exponential growth as shown in Figure~\ref{fig:voc}.

The latest analysis from the ONS CIS for the B.1.1.7 variant is shown in
Figure~\ref{fig:ONS} and shows dynamics that are arguably consistent with
a `stochastic delay then exponential takeoff' model of strain emergence.

\section{Dynamic importation and static control}

We start by revisiting the simple analytic results of \citet{ScaliaTomba:2008}.
These were derived in response to simulations relating to border controls
(e.g.\ \citet{Hollingsworth:2006}). We will derive these slightly differently
from the original paper, and also generalise.

First, let $X(t)$ be the cumulative number of importations at time $t$ and let
$q$ be the probability that each introduction eventually becomes extinct (we
assume these probabilities are independent). Then the probability that there
has not yet been an introduction of a new variant that will lead to a large
outbreak is
\be
\sum_{z=0}^{\infty} \mathrm{Pr}(X(t) = x) q^x = G_{X(t)}(q) \text{ ,}
\ee
where $G$ stands for a pgf -- using the expression for a Poisson random
variable with mean $\Lambda(t)$ (although we could generalise to a negative
binomial because infections tend to be clustered as seen in the data) we recover the main
result for the time $T$ of failure to keep the variant out
from \citeauthor{ScaliaTomba:2008} of
\be
F_T(t) = \mathrm{Pr}(T \leq t) = 1 - \exp(-(1-q)\Lambda(t)) \text{ .}
\label{wst}
\ee
We use this equation for the choice $\Lambda(t) = 1.7 t$ to produce
Figure~\ref{fig:F}.  As \citeauthor{ScaliaTomba:2008} note, substituting
\be
\Lambda(t) = \Lambda_0 {\rm e}^{rt}
\text{ ,} \qquad
t_D = \frac{\log(2)}{r} \text{ ,}
\ee
into \eqref{wst}, we get the following expression for the difference in median
times of introduction under control of
\be
\tau_q - \tau_0 = \frac{-\log(1-q)}{\log(2)} t_D \text{ .}
\ee
Calling this delay in median time $\delta \tau$, we can derive an expression for
the level of control needed to achieve this delay:
\be
q(\delta \tau) = 1- \exp\left( - \frac{\delta \tau}{t_D} \log(2)\right) \text{ .}
\ee
Evaluating this expression for some realistic parameter values gives
Figure~\ref{fig:q}.

\section{Dynamic importation and onward transmission}

Now we consider disease dynamics -- this is done in a single-type / SIR
framework for simplicity, but an equivalent multitype case is possible
\citep{Dorman:2004}. In this model, each particle produces a number of
offspring that is an iid copy of random variable $Y$ after a time that is
$\mathrm{Exp}(\beta + \gamma)$ distributed, where
\be
G_{Y}(s)= \sum_{y=0}^{\infty} \mathrm{Pr}(Y = y) s^y = 
\frac{1}{\beta + \gamma} \left(\beta\big(1+\frac{1}{k}(1-s)\big)^{-k}s + \gamma\right)
 \text{ .} \label{SSEp}
\ee
In order to allow superspreading events (SSEs) to occur, we have chosen this
offspring distribution so that the number of newly infected cases is drawn from
a negative binomial distribution with mean $\frac{\beta}{\beta+\gamma}$ and
dispersion parameter $k$ \citep{Lloyd-Smith:2005}. For $k=-1$ we recover
standard infection dynamics, with \citet{Endo:2020} estimating $k\approx 0.1$.

Using this offspring distribution, we define $q(t)$ as the probability that the
lineage of a single imported case has become extinct by time $t$ (generalising
the static $q$ above). Note that, in both cases, the mean number of secondary
cases from a single infection event is given by $\frac{\beta}{\beta+\gamma} =
\left.\frac{{\mathrm{d}}P}{{\mathrm{d}}s}\right|_{s=1}$. As noted in
\citet{Dorman:2004} is equivalent to a birth-death chain in which the expected
number of secondary cases produced over by a single individual over the course
of their infection is $\frac{\beta}{\gamma} = R$.

Now let $Z(t)$ be the number of particles in a branching process model with
importation of particles at rate $\lambda(t)$ with integrated intensity 
\be
\Lambda(t) = \int_{u=0}^{t} \lambda(u) \mathrm{d}u
\ee
We can then generalise the results of \citeauthor{ScaliaTomba:2008} to obtain
the probability that there is a positive number of cases time at time $t$,
\be
\mathrm{Pr}(Z(t) = 0) = 1 - {\mathrm{exp}}(-(1 - q(t))\Lambda(t))
\label{Z0}
\ee
Since this quantity is very close to zero for most realistic scenarios, it is
most natural to investigate the behaviour of $q(t)$ through solving the
Kolmogorov equations for the branching process for different values of $k$ and
reproduction number $R$ for the new variant in the UK environment. These
results are shown in Figure~\ref{fig:q_t}.

\bibliography{mybib}
\bibliographystyle{abbrvnat}

\clearpage

\begin{figure}
\centering
\includegraphics[width=1.0\textwidth]{./voc_c.pdf}\\
\includegraphics[width=1.0\textwidth]{./voc_la.pdf}\\
\includegraphics[width=1.0\textwidth]{./voc_r.pdf}
\caption{Analysis of data from the variant of concern.}
\label{fig:voc}
\end{figure}

\clearpage

\begin{figure}
\centering
\fbox{
\includegraphics[width=0.99\textwidth]{./ons_sgtf.pdf}
}
\caption{\textbf{These outputs have been provided as management information for
operational planning purposes. They are provisional estimates circulated before
public release. This management information should not be shared widely and
should only be used for operational planning purposes. This information is not
to be used publicly and ahead of being released publicly by ONS.} Prevalence of
wildtype and B.1.1.7 in the UK provided by the ONS CIS team, showing relatively
long periods in some geographies before the establishment of exponential
growth.}
\label{fig:ONS}
\end{figure}

\clearpage

\begin{figure}
\centering
\includegraphics[width=0.9\textwidth]{./var_F.pdf}
\caption{Constant importation rate: probability of established introduction over
time for different probabilities of stopping an introduction.}
\label{fig:F}
\end{figure}

\begin{figure}
\centering
\includegraphics[width=0.9\textwidth]{./var_q.pdf}\\
\includegraphics[width=0.9\textwidth]{./var_q_log.pdf}
\caption{Exponentially growing importation rate: Levels of control needed to shift median
introduction time by a desired amount.}
\label{fig:q}
\end{figure}

\clearpage

\begin{figure}
\centering
\includegraphics[width=1.0\textwidth]{./q_t_k-1.0.pdf}\\
\includegraphics[width=1.0\textwidth]{./q_t_k0.1.pdf}\\
\includegraphics[width=1.0\textwidth]{./q_t_k1.0.pdf}
\caption{Extinction probabilities for single lineages over time.}
\label{fig:q_t}
\end{figure}



\end{document}



